% -------------------------------------------------------------------------
% WBS ID: RES-MOD-SRC
% Deliverable: Physical Source Terms and Constitutive Closures
% Status: Draft complete
% Evidence status: Research specification complete for regional NLSWE sources and local URANS baseline closure
% Review status: Not reviewed
% Last updated: 2026-07-18
% Dependencies: RES-MOD-EQS, RES-PHY-SRC, RES-DAT-GEO, RES-DAT-SRC
% Downstream interfaces: RES-NUM-SRC, RES-NUM-TIM, RES-VER-SPEC
% -------------------------------------------------------------------------

\section{RES-MOD-SRC --- Physical Source Terms and Constitutive Closures}
\label{sec:res-mod-src}

\subsection{Purpose and Scope}
\label{sec:res-mod-src-purpose}

% TODO: Define what this Deliverable covers and explicitly excludes.

\subsection{Research Questions}
\label{sec:res-mod-src-research-questions}

% TODO: State the questions that must be answered before the Deliverable
% can be accepted.

\subsection{Terminology and Definitions}
\label{sec:res-mod-src-definitions}

% TODO: Define the technical terminology, symbols, quantities and
% classifications used in this Deliverable.

\subsection{Literature Evidence}
\label{sec:res-mod-src-literature-evidence}

% TODO: Record evidence from peer-reviewed papers, authoritative datasets,
% standards, technical reports and primary documentation.
%
% Do not create a detached literature-summary list. Explain how each source
% supports, contradicts or limits a project decision.

\subsubsection{Kozelkov et al. (2022)}
\label{sec:res-mod-src-evidence-kozelkov-2022}

\textcite{kozelkov2022turbulence} compared a two-phase
Navier--Stokes calculation without a turbulence closure against a
Reynolds-averaged calculation using the $k$--$\omega$ SST model.
Differences were small during smooth propagation but became material
during breaking, run-up, overtopping, collapse, and barrier loading.
The reported changes in total barrier force reached engineering
significance in the more energetic cases.

\subsubsection{Watanabe et al. (2022)}
\label{sec:res-mod-src-evidence-watanabe-2022}

\textcite{watanabe2022validation} compared blind submissions using
laminar, Reynolds-averaged, and large-eddy formulations. Dynamic
sub-grid kinetic-energy LES submissions generally produced more
consistent force predictions than fixed-coefficient Smagorinsky
submissions, while the submitted $k$--$\varepsilon$ calculations were
more stable and consistent than the submitted $k$--$\omega$
calculations for that experiment.

\subsubsection{Menter (1994)}
\label{sec:res-mod-src-evidence-menter-1994}

\textcite{menter1994twoequation} introduced the $k$--$\omega$ SST
eddy-viscosity closure by blending near-wall $k$--$\omega$ behaviour
with outer-flow $k$--$\varepsilon$ behaviour and limiting turbulent
shear stress. Source role: `3D-SST`. Project conclusion: SST is a
defensible operational URANS baseline for high-Reynolds-number local
tsunami run-up and impact studies. Limitation: this source is a
general turbulence-model source; project validation must still test the
closure against free-surface, pressure, overtopping and load outputs.

\subsection{Turbulence Treatment for the Local Three-Dimensional Model}
\label{sec:res-mod-src-local-turbulence}

The selected parent Navier--Stokes equations require a treatment of
unresolved turbulent transport in the local impact region.

The local regime is treated as high-Reynolds-number turbulent
free-surface impact flow. A molecular-viscosity-only calculation may be
used diagnostically but is not accepted as the operational model for
candidate screening.

These results do not establish one universally optimal closure. The
submitted models differed in their meshes, solvers, interface
treatments, wall conditions, and numerical schemes. The findings
instead establish that:

\begin{enumerate}
  \item turbulence closure is required in the local breaking and
    impact domain;
  \item the closure must be assessed within a controlled common
    numerical framework;
  \item free-surface agreement alone is insufficient for selecting
    the closure;
  \item closure selection must include pressure, force, impulse,
    overtopping, and downstream-flow measures.
\end{enumerate}

\subsubsection{Turbulence-Model Hierarchy}
\label{sec:res-mod-src-local-turbulence-hierarchy}

The turbulence formulations are assigned the following project roles:

\begin{table}[htbp]
  \centering
  \small
  \renewcommand{\arraystretch}{1.25}
  \begin{tabular}{
      p{0.22\textwidth}
      p{0.27\textwidth}
      p{0.40\textwidth}
    }
    \hline
    Formulation
    &
    Project role
    &
    Current decision
    \\
    \hline
    Unsteady RANS
    &
    Operational local-impact model
    &
    Selected as the initial candidate-evaluation formulation with
    $k$--$\omega$ SST as the baseline closure
    \\
    Dynamic sub-grid kinetic-energy LES
    &
    High-fidelity comparison model
    &
    Retained for validation cases and shortlisted geometries
    \\
    Hybrid RANS--LES
    &
    Possible intermediate-fidelity model
    &
    Retained for later research; insufficient evidence for current
    selection
    \\
    Fixed-coefficient Smagorinsky LES
    &
    Alternative LES closure
    &
    Retained as a comparator but not provisionally preferred
    \\
    Molecular-viscosity-only Navier--Stokes
    &
    Diagnostic case
    &
    Not accepted as the operational high-Reynolds-number impact
    model
    \\
    \hline
  \end{tabular}
  \caption{Project roles assigned to the candidate turbulence
  formulations}
  \label{tab:res-mod-src-local-turbulence-hierarchy}
\end{table}

The operational and high-fidelity formulations shall use the same
parent fluid equations, material properties, geometry, incident state,
and output definitions wherever practical.

\subsection{Governing Relations and Quantitative Definitions}
\label{sec:res-mod-src-governing-relations}

% TODO: Record relevant equations, dimensionless groups, empirical models,
% extraction rules or quantitative definitions.
%
% Use align environments for displayed equations.
% Every displayed equation must have a unique \label{eq:...}.
% Do not place punctuation after displayed equations.

\subsubsection{Unsteady Reynolds-Averaged Formulation}
\label{sec:res-mod-src-local-urans}

The instantaneous fluid velocity is decomposed as:

\begin{align}
  u_i
  &=
  \overline{u}_i
  +
  u_i^{\prime}
  \label{eq:res-mod-src-urans-velocity-decomposition}
\end{align}

Averaging introduces the Reynolds stress:

\begin{align}
  \tau_{ij}^{\mathrm{R}}
  &=
  -\rho_{\mathrm{f}}
  \overline{
    u_i^{\prime}u_j^{\prime}
  }
  \label{eq:res-mod-src-urans-reynolds-stress}
\end{align}

For a linear eddy-viscosity closure:

\begin{align}
  \tau_{ij}^{\mathrm{R}}
  &=
  2\mu_{\mathrm{t}}
  \overline{S}_{ij}
  -
  \frac{2}{3}
  \rho_{\mathrm{f}}k\delta_{ij}
  \label{eq:res-mod-src-urans-boussinesq-closure}
  \\
  \overline{S}_{ij}
  &=
  \frac{1}{2}
  \left(
    \frac{\partial\overline{u}_i}{\partial x_j}
    +
    \frac{\partial\overline{u}_j}{\partial x_i}
  \right)
  \label{eq:res-mod-src-urans-strain-rate}
\end{align}

The operational URANS closure is the $k$--$\omega$ SST model. Its
generic transport form is recorded as:

\begin{align}
  \frac{\partial\left(\rho_{\mathrm{f}} k\right)}{\partial t}
  +
  \nabla\cdot
  \left(
    \rho_{\mathrm{f}}\mathbf{u}_{\mathrm{f}} k
  \right)
  &=
  \nabla\cdot
  \left[
    \left(
      \mu_{\mathrm{f}}
      +
      \sigma_k\mu_{\mathrm{t}}
    \right)
    \nabla k
  \right]
  +
  P_k
  -
  \beta^{*}\rho_{\mathrm{f}}k\omega
  \label{eq:res-mod-src-sst-k-equation}
  \\
  \frac{\partial\left(\rho_{\mathrm{f}}\omega\right)}{\partial t}
  +
  \nabla\cdot
  \left(
    \rho_{\mathrm{f}}\mathbf{u}_{\mathrm{f}}\omega
  \right)
  &=
  \nabla\cdot
  \left[
    \left(
      \mu_{\mathrm{f}}
      +
      \sigma_{\omega}\mu_{\mathrm{t}}
    \right)
    \nabla\omega
  \right]
  +
  \gamma
  \frac{\omega}{k}
  P_k
  -
  \beta\rho_{\mathrm{f}}\omega^2
  +
  D_{\omega}
  \label{eq:res-mod-src-sst-omega-equation}
\end{align}

The eddy viscosity is limited in SST form by:

\begin{align}
  \mu_{\mathrm{t}}
  &=
  \rho_{\mathrm{f}}
  \frac{
    a_1 k
  }{
    \max
    \left(
      a_1\omega,\,
      S F_2
    \right)
  }
  \label{eq:res-mod-src-sst-eddy-viscosity}
\end{align}

where $P_k$ is the turbulence-production term, $S$ is a strain-rate
measure, $F_2$ is the SST blending function and $D_{\omega}$ contains
the cross-diffusion contribution. The coefficient set and blending
functions shall follow the selected SST reference rather than being
retuned during the research specification.

Decision trace: source role `3D-SST`; supporting sources
\textcite{menter1994twoequation}, \textcite{kozelkov2022turbulence}
and \textcite{watanabe2022validation}. Project-specific conclusion:
URANS $k$--$\omega$ SST is the baseline turbulence closure for repeated
local candidate assessment. Limitation: the closure is not accepted
for final engineering certification until the validation hierarchy in
`RES-VER` confirms pressure, force, impulse and overtopping behaviour.
Downstream handoff: `RES-NUM-SPEC` shall define when $k$ and $\omega$
are solved in the time-step sequence and `RES-VER-MET` shall define
the URANS--LES comparison metrics.

\subsubsection{Large-Eddy Formulation}
\label{sec:res-mod-src-local-les}

Large-eddy simulation filters the velocity field into resolved and
sub-grid components:

\begin{align}
  u_i
  &=
  \widetilde{u}_i
  +
  u_i^{\mathrm{sgs}}
  \label{eq:res-mod-src-les-velocity-decomposition}
\end{align}

Filtering introduces the sub-grid stress:

\begin{align}
  \tau_{ij}^{\mathrm{sgs}}
  &=
  \rho_{\mathrm{f}}
  \left(
    \widetilde{u_i u_j}
    -
    \widetilde{u}_i\widetilde{u}_j
  \right)
  \label{eq:res-mod-src-les-subgrid-stress}
\end{align}

The deviatoric sub-grid stress may be represented as:

\begin{align}
  \tau_{ij}^{\mathrm{sgs}}
  -
  \frac{1}{3}
  \tau_{kk}^{\mathrm{sgs}}
  \delta_{ij}
  &=
  -2\mu_{\mathrm{sgs}}
  \widetilde{S}_{ij}
  \label{eq:res-mod-src-les-eddy-viscosity-closure}
\end{align}

A dynamic sub-grid kinetic-energy model is retained as the provisional
high-fidelity candidate. Its exact transport equation and coefficient
evaluation shall be selected during the dedicated turbulence-closure
study.

\subsubsection{Turbulence-Closure Selection Gate}
\label{sec:res-mod-src-local-closure-gate}

The operational closure shall be accepted only after comparison against
the high-fidelity formulation for representative impact cases.

The comparison shall include:

\begin{itemize}
  \item incident and transmitted free-surface histories;
  \item pressure distributions;
  \item resultant forces and moments;
  \item force impulse;
  \item overtopping volume;
  \item reflected and transmitted energy;
  \item downstream recirculation and wake structure;
  \item resolved and modelled turbulent kinetic energy;
  \item sensitivity of the barrier-performance ranking.
\end{itemize}

The operational URANS model may be retained where it preserves the
design ranking and the required engineering quantities within the
acceptance tolerances defined under `RES-VER-MET`.

If the barrier ranking changes materially between URANS and LES, the
lower-fidelity model shall not be used as the sole optimisation
evaluator until the source of the discrepancy has been resolved.

\subsection{System Boundary}
\label{sec:res-mod-src-system-boundary}

% TODO: Complete this RES-MOD Domain-specific subsection.

\subsection{State Variables and Parameters}
\label{sec:res-mod-src-state-variables-parameters}

% TODO: Complete this RES-MOD Domain-specific subsection.

\subsection{Continuous Governing Model}
\label{sec:res-mod-src-continuous-governing-model}

% TODO: Complete this RES-MOD Domain-specific subsection.

\subsection{Source Terms and Closures}
\label{sec:res-mod-src-source-terms-closures}

The regional two-dimensional NLSWE source catalogue is established at
research-specification level. The baseline retained terms are bed slope,
Manning bottom friction, dynamic earthquake bathymetry, and optional
sponge damping near artificial open boundaries. Coriolis and lateral
eddy viscosity are not part of the baseline regional model; they remain
optional sensitivity terms if the selected domain scale or validation
evidence requires them.

The continuous regional source split is:

\begin{align}
  \mathbf{S}
  &=
  \mathbf{S}_{\mathrm{b}}
  +
  \mathbf{S}_{\mathrm{f}}
  +
  \mathbf{S}_{\mathrm{sp}}
  \label{eq:res-mod-src-regional-source-split}
\end{align}

where $\mathbf{S}_{\mathrm{b}}$ is the bed-slope source,
$\mathbf{S}_{\mathrm{f}}$ is Manning bottom friction, and
$\mathbf{S}_{\mathrm{sp}}$ is the sponge-layer relaxation source
defined in \cref{sec:res-mod-eqs-continuous-governing-model}. The
earthquake contribution is applied by updating the bed field
$b(x,y,t)$; no additional depth-equation mass source is added in the
conserved-depth formulation.

\subsection{Initial, Boundary and Interface Conditions}
\label{sec:res-mod-src-initial-boundary-interface-conditions}

% TODO: Complete this RES-MOD Domain-specific subsection.

\subsection{Model Coupling Requirements}
\label{sec:res-mod-src-model-coupling-requirements}

% TODO: Complete this RES-MOD Domain-specific subsection.

\subsection{Sufficient Conditions for Validity}
\label{sec:res-mod-src-sufficient-conditions-validity}

% TODO: Complete this RES-MOD Domain-specific subsection.

\subsection{Candidate Approaches}
\label{sec:res-mod-src-candidate-approaches}

% TODO: Define the credible candidate approaches considered.

\subsection{Critical Comparison}
\label{sec:res-mod-src-critical-comparison}

% TODO: Compare candidates using physically and project-relevant criteria.
% Do not select an approach solely on popularity or implementation ease.

\subsection{Assumptions and Applicability}
\label{sec:res-mod-src-assumptions}

% TODO: State assumptions and the conditions under which they are valid.

\subsection{Limitations and Failure Conditions}
\label{sec:res-mod-src-limitations}

% TODO: State where the evidence, model or selected approach ceases to be
% reliable or applicable.

\subsection{Project-Specific Conclusions}
\label{sec:res-mod-src-conclusions}

% TODO: Convert the evidence into conclusions specific to the Studentship
% project. Do not merely restate literature findings.

For the regional model, the source-term research gate is closed at
specification level. Remaining work is numerical verification of the
well-balanced bed-source treatment, validation of the selected
earthquake source against observations, and selection of data values for
roughness, bathymetry, and source timing. These are not changes to the
regional mathematical model unless the validation evidence rejects the
specified formulation.

\subsection{Selected Approach and Rationale}
\label{sec:res-mod-src-selected-approach}

% TODO: Record the selected approach only after sufficient evidence exists.
% State the alternatives rejected and the reasons for rejection.

The selected regional source approach is:

\begin{itemize}
  \item bed-slope forcing with hydrostatic reconstruction and matching
    bed-source discretisation;
  \item Manning bottom friction, advanced by a split semi-implicit
    update in the numerical method;
  \item dynamic moving-bed earthquake source,
    $b=b_0+\sum_m\Delta b_m R_m(t)$;
  \item sponge-layer relaxation at artificial ocean boundaries where
    required by domain truncation;
  \item Coriolis and explicit lateral eddy viscosity disabled in the
    baseline, retained only as sensitivity options.
\end{itemize}

The selected local closure approach is:

\begin{itemize}
  \item high-Reynolds-number URANS as the operational local
    hydrodynamic model;
  \item $k$--$\omega$ SST as the baseline eddy-viscosity closure;
  \item LES retained for selected high-fidelity comparison cases;
  \item molecular-viscosity-only Navier--Stokes retained only as a
    diagnostic calculation;
  \item porous, compressible-air, sediment, scour and structural
    damage source terms deferred until their activation gates are met.
\end{itemize}

\subsection{Unresolved Decisions}
\label{sec:res-mod-src-unresolved-decisions}

% TODO: Record unresolved questions, required evidence and decision owners.

The mathematical form of the regional source terms and the local URANS
baseline closure are no longer unresolved. The open items are data,
verification and lower-level closure-implementation items:

\begin{itemize}
  \item final Manning coefficient fields;
  \item final static bathymetry/topography dataset;
  \item final fault-source dataset and rupture-time functions;
  \item lake-at-rest and moving-bottom verification evidence;
  \item decision on whether optional Coriolis sensitivity is required;
  \item local wall-function constants and roughness treatment for the
    selected case-study materials;
  \item final LES sub-grid closure and comparison-case set.
\end{itemize}

\subsection{Requirements and Handoffs}
\label{sec:res-mod-src-requirements}

% TODO: State requirements passed to other Research Domains, Software,
% verification activities or later execution work.

`RES-NUM-SRC` shall implement the source terms consistently with the
finite-volume face fluxes, preserve still-water balance over non-flat
bathymetry, and report any loss of positivity at wet/dry fronts.
`RES-DAT-GEO` and `RES-DAT-SRC` shall provide the static bed,
coordinate/datum transform, and fault-source data required by the
source definitions.

`RES-NUM-SPEC` shall discretise
\cref{eq:res-mod-src-sst-k-equation,eq:res-mod-src-sst-omega-equation}
after the pressure-corrected velocity and face flux have been accepted
for the time step. `RES-PHY-SCL` shall justify the high-Reynolds-number
regime, and `RES-VER-MET` shall prevent free-surface agreement alone
from being treated as closure validation.

\subsection{Deliverable Completion Record}
\label{sec:res-mod-src-completion-record}

% TODO: Record whether the Deliverable is:
%
% - Not started
% - In progress
% - Draft complete
% - Reviewed
% - Accepted
%
% Acceptance requires:
%
% - the research questions to be answered;
% - claims to be supported by appropriate evidence;
% - assumptions and limitations to be explicit;
% - the project-specific conclusion to be stated;
% - downstream requirements to be traceable;
% - unresolved decisions to be closed or formally deferred.
